\chapter{isudbvisuf}

\cite{Demmel1983}
\begin{definition} % Wilkinson
    Let $T \in \C^{d\times d}$. We say that $T$ is defective if its Jordan canonical form is not strictly diagonal, i.e., if there exists a non-trivial Jordan block (of dimension $>1$). Equivalently, $T$ has fewer than $d$ linearly independent eigenvectors.
\end{definition}
\begin{definition}
    For a matrix $T \in \C^{d \times d}$, we define the distance to the nearest defective matrix as:
    $\delta(T) := \inf \{\,\|A - T\| \, : \, A \in \C^{d \times d} \text{ and } A \text{ is defective }\}$.
\end{definition}

\begin{theorem}[Schur]
    Let $T\in\C^{d\times d}$ and let $\sigma(T) = \left(\lambda_1, \dots, \lambda_d\right)$ denote its spectrum (with multiplicities counted) in an arbitrary order. Then there exists a unitary matrix $Q$ and an upper-triangular matrix $U$, such that $T = Q^*UQ$ and $\mathrm{diag}(U) = \left(\lambda_1, \dots, \lambda_d\right)$.
\end{theorem}
\begin{proof}
    See \cite[p. 79]{HornJohnson1985}.
\end{proof}


\begin{lemma}\label{lemma} % Demmel lemma 2.1
    Let $\| \cdot \|$ denote either the Frobenius norm or the operator 2-norm on $\C^{d\times d}$. Let $P$ be a property depending only on Jordan canonical form of a matrix. For $T_0 \in \C^{d\times d}$ and a non-singular $S \in \C^{d \times d}$ define $\delta := \inf_{T \text{ satisfies } P}  \|T_0 - T\|$ and $\delta_S := \inf_{T \text{ satisfies } P} \|S^{-1}T_0S - T\|$. Then,
    \begin{equation*}
        \frac{\delta}{\kappa(S)} \leqslant \delta_S \leqslant \delta\kappa(S).
    \end{equation*}
\end{lemma}
\begin{proof}
    The property $P$ depends only on the Jordan canonical form of the matrix, therefore $T$ satisfies $P$ if and only if $S^{-1}TS$ satisfies $P$ for arbitrary non-singular matrix $S$. Using \eqref{eq:matrix-norm-inequality} and the definition of $\kappa(S)$, we get:
    \begin{equation*}
        \left\|S^{-1} T_0 S - S^{-1} T S \right\| \leqslant \left\|S^{-1}\right\| \left\|T_0 - T\right\| \left\|S\right\| = \|T_0-T\|\kappa(S).
    \end{equation*}
    This proofs the second equality. The first one can be shown analogously.
\end{proof}
\begin{theorem}
Let $T \in \C^{d \times d}$ and $\sigma(T) = (\lambda_1, \dots, \lambda_d)$ be the spectrum of $T$ (with multiplicities counted). Then: 
\begin{equation*}
    \delta(T) \leqslant \min_{i\neq j} |\lambda_i - \lambda_j|.
\end{equation*}
\end{theorem}
\begin{proof}
    Without loss of generality, assume $\min_{i\neq j} |\lambda_i - \lambda_j| = |\lambda_1 - \lambda_2| =: \varepsilon$. We will show that $\delta(T) \leqslant \varepsilon$. 
    
    Let $T = Q^T U Q$ be the Schur decomposition of $T$ such that $\mathrm{diag}(U) = (\lambda_1, \lambda_2, \dots, \lambda_d)$. Define $U' = \left(U'_{i,j}\right)_{i,j=1}^d$ by: 
    \begin{equation*} 
        U'_{i,j} = \begin{cases}
            \lambda_1, \quad \text{if } i=j=2 \\
            U_{i,j}, \quad \text{else}
        \end{cases}.
    \end{equation*}
   In other words, $U'$ is obtained from $U$ by replacing $\lambda_2$ with $\lambda_1$ on the diagonal. Set $T' := Q^* U' Q$. Since $Q$ is unitary and $\kappa(Q) = 1$, Lemma \ref{lemma} gives $\| T - T' \| \leqslant \kappa(Q) \cdot \|U - U'\| \leqslant \varepsilon $.
    
     The matrix $T'$ has a multiple eigenvalue $\lambda_1$. Let $T' = S^{-1} J S$ be its Jordan canonical form. If $T'$ is already defective at $\lambda_1$, we are done. Otherwise, $J$ contains at least two $1 \times 1$ Jordan blocks corresponding to eigenvalue $\lambda_1$. We choose $\eta > 0$ and construct $J'$ by replacing those two blocks with a non-trivial $2\times2$ Jordan block with $\eta$ on the superdiagonal:
     \begin{equation*}
         J' = \begin{pmatrix}
         \ddots & & & \\
         & \lambda_1 & \eta & \\
         & & \lambda_1 & \\
         & & & \ddots
         \end{pmatrix}.
     \end{equation*}
     Let $T'':=S^{-1}J'S$. By Lemma \ref{lemma}, $\|T' - T''\| \leqslant \|J - J'\|\cdot\kappa(S) \leqslant \eta \cdot\kappa(S)$, and by construction, $T''$ is defective.

     Thus, we have obtained a defective matrix in a $(\varepsilon + \eta \cdot \kappa(S))$-neighbourhood of $T$. Since $\eta$ was arbitrary, the desired inequality follows.
    
\end{proof}